\section{High-Dimensional Methods}
\red{needs refinement. The entire chapter is not curated yet. Notes are rough, need better interpretation and intuition (e.g. lasso soft threshold + coordinate descent as residual compensation to check for normal orthogonality in linear regression $\E[Xe]=0$ with $\beta_{-j}$ fixed to decide whether to include $\beta_j$). log-likelihood + LQA + gradient descent procedures.}

High-dimensional regression uses penalization to make estimation, prediction, and variable selection feasible when $p$ is large relative to $n$. (See Notes pp.~1--3.)

\begin{block}
\textbf{Big-picture map of the chapter.}
\begin{enumerate}
    \item The chapter starts with ridge regression: an $L_2$ penalty stabilizes least squares, gives a unique estimator even when $X^\top X$ is singular, and acts as shrinkage rather than variable selection. (See Notes pp.~4--15.)
    \item It then introduces Lasso and related $L_q$ penalties. The central point is the tension between convexity, sparsity, and bias: the $L_1$ penalty is the convex penalty that can set coefficients exactly to zero. (See Notes pp.~16--28.)
    \item The notes next turn to computation. Coordinate gradient descent solves the linear-regression Lasso through repeated soft-thresholding, and subgradient / KKT conditions give the algebraic certificate for optimality. (See Notes pp.~27, 32--40, 42--53.)
    \item After that, the chapter proves a non-asymptotic model-selection consistency result for Lasso under sparsity, beta-min, and irrepresentable-design assumptions. (See Notes pp.~41, 48--61.)
    \item The same penalization logic is then extended to generalized linear models, where local quadratic approximation, MM, and proximal gradient descent turn smooth loss plus nonsmooth penalty problems into repeated easier subproblems. (See Notes pp.~62--78.)
    \item The final part studies alternatives to ordinary Lasso: SCAD and adaptive Lasso aim at oracle behavior by reducing large-signal bias, while elastic net adds an $L_2$ component to encourage grouping among highly correlated predictors. (See Notes pp.~79--115.)
\end{enumerate}
\end{block}

\begin{mybox}
\textbf{Supplementary problem-solving hierarchy for high-dimensional regression.}
\begin{enumerate}
    \item First identify the goal: prediction, variable selection, or inference about the selected support. Ridge is mainly for prediction/stabilization; Lasso is for sparse prediction and selection; SCAD/adaptive Lasso target oracle-style support recovery; elastic net targets correlated groups. (See Notes pp.~15--21, 79--90, 103--113.)
    \item Then identify the optimization structure: smooth convex loss plus $L_2$ penalty gives closed-form ridge; smooth convex loss plus $L_1$ penalty gives coordinate descent or proximal gradient; nonconvex penalties such as SCAD require local-optimizer thinking. (See Notes pp.~6, 32--40, 68--77, 87--89.)
    \item For theory questions, separate the statistical assumptions from the algorithm. Lasso selection consistency uses KKT conditions plus concentration; SCAD oracle theory is about the existence of a good local maximizer, not automatic global optimization. (See Notes pp.~49--61, 87--102.)
\end{enumerate}
\end{mybox}

\subsection{Ridge regression and shrinkage}
\subsubsection{Core setup}
For centered $y$ and standardized columns of $X$, ridge regression solves (See Notes pp.~4--6.)
\[
\hat\beta^{\mathrm{ridge}}
= \argmin_{\beta}
\left\{\norm{y-X\beta}_2^2 + \lambda \norm{\beta}_2^2\right\}.
\]
Equivalently, it solves least squares subject to $\norm{\beta}_2^2 \le t$, with $\lambda$ as the Lagrange multiplier. The intercept is not penalized; after centering $y$, one can take $\beta_0=0$ in the matrix form. (See Notes pp.~4--6.)

The ridge solution is
\[
\hat\beta^{\mathrm{ridge}}
=(X^\top X+\lambda I)^{-1}X^\top y.
\]
The stabilizing effect is the main reason ridge matters in high dimensions: even if $p>n$ and $X^\top X$ is singular, $X^\top X+\lambda I$ is positive definite for $\lambda>0$. (See Notes p.~6.)

\subsubsection{Shrinkage, bias, and Bayesian interpretation}
\red{[summarize and clean up]}
Let $M=X^\top X$ and assume fixed design. The notes show (See Notes pp.~8--9.)
\[
\hat\beta^{\mathrm{ridge}}
= (I+\lambda M^{-1})^{-1}\hat\beta^{\mathrm{ls}},
\qquad
\E(\hat\beta^{\mathrm{ridge}})
= (I+\lambda M^{-1})^{-1}\beta,
\]
so ridge is biased unless $\lambda=0$. Under an orthonormal design ($M=I_p$), this becomes the scalar shrinkage formula
\[
\hat\beta^{\mathrm{ridge}}_j
= \frac{1}{1+\lambda}\hat\beta^{\mathrm{ls}}_j,
\]
up to the scaling convention used for $X$. (See Notes pp.~6, 38.)

The Bayesian interpretation is: if $y_i\mid x_i,\beta \sim N(\beta_0+x_i^\top\beta,\sigma^2)$ and $\beta\sim N(0,\tau^2I)$, then the negative log posterior is proportional to
\[
\norm{y-X\beta}_2^2 + \frac{\sigma^2}{\tau^2}\norm{\beta}_2^2,
\]
so ridge corresponds to a Gaussian prior with $\lambda=\sigma^2/\tau^2$ under the notes' unscaled objective. (See Notes p.~9.)

\subsubsection{SVD computation and effective degrees of freedom}
Using an SVD $X=UDV^\top$ with singular values $d_1\ge \cdots\ge d_p$, ridge can be written as (See Notes pp.~10--12.)
\[
\hat\beta^{\mathrm{ridge}}
=\sum_{j=1}^p
\left(\frac{d_j}{d_j^2+\lambda}u_j^\top y\right)v_j,
\qquad
\hat y^{\mathrm{ridge}}
=\sum_{j=1}^p
\left(\frac{d_j^2}{d_j^2+\lambda}u_j^\top y\right)u_j.
\]
Thus ridge first projects $y$ onto the left singular-vector directions and then shrinks the coordinates by $d_j^2/(d_j^2+\lambda)$.

The ridge smoother matrix is
\[
S_\lambda=X(X^\top X+\lambda I)^{-1}X^\top,
\]
and its effective degrees of freedom are
\[
\operatorname{df}(\lambda)=\tr(S_\lambda)
= \sum_{j=1}^p \frac{d_j^2}{d_j^2+\lambda}.
\]
This decreases with $\lambda$ and equals $\operatorname{rank}(X)$ when $\lambda=0$. (See Notes pp.~13--15.)

\begin{mybox}
\textbf{What to remember about ridge.}
\begin{itemize}
    \item \underline{Existence and uniqueness:} ridge fixes the singularity of $X^\top X$ by adding $\lambda I$. (See Notes pp.~6, 15.)
    \item \underline{Shrinkage but no selection:} ridge coefficients are pulled toward zero but are generally not exactly zero. (See Notes pp.~6, 15, 17.)
    \item \underline{Useful diagnostics:} SVD explains the shrinkage direction by direction, and $\operatorname{df}(\lambda)=\tr(S_\lambda)$ measures the effective model complexity. (See Notes pp.~10--15.)
\end{itemize}
\end{mybox}

\subsection{Lasso, tuning, and variable-selection targets}
\subsubsection{Core objects}
\red{[make this align with ridge objective]}
The Lasso solves (See Notes pp.~16--18.)
\[
\hat\beta^{\mathrm{lasso}}
= \argmin_{\beta}
\left\{\sum_{i=1}^N
(y_i-\beta_0-x_i^\top\beta)^2
+\lambda\sum_{j=1}^p \abs{\beta_j}\right\},
\]
or equivalently least squares subject to $\sum_j \abs{\beta_j}\le t$. The singularity of the $L_1$ penalty at zero is what allows exact zeros in $\hat\beta$, so Lasso can perform variable selection while ridge cannot. (See Notes pp.~16--18.)

The notes place ridge and Lasso inside the larger class
\[
\argmin_\beta
\left\{
\sum_{i=1}^N(y_i-\beta_0-x_i^\top\beta)^2
+\lambda\sum_{j=1}^p \abs{\beta_j}^q
\right\}.
\]
Here $q=0$ corresponds to best subset selection, $q=1$ to Lasso, and $q=2$ to ridge. The $L_q$ penalty is singular at zero exactly when $q\le 1$, while the objective is convex exactly when $q\ge 1$. Therefore $q=1$ is the boundary case that gives both convex optimization and sparsity. (See Notes pp.~18--19.)

\red{add plot with annotations}

\subsubsection{Tuning parameter selection}
The tuning parameter $\lambda$ controls the amount of regularization: large $\lambda$ shrinks more aggressively, and for Lasso large enough $\lambda$ sets all coefficients to zero. (See Notes pp.~21--23, 39.)

\red{remove equation. it's just MSE on validation set}
For $K$-fold cross validation, partition the training data into folds $T_1,\ldots,T_K$, fit $\hat\beta_\lambda^{(-k)}$ without fold $k$, compute
\[
\operatorname{err}_k(\lambda)
=\abs{T_k}^{-1}\sum_{(x_i,y_i)\in T_k}
(y_i-x_i^\top\hat\beta_\lambda^{(-k)})^2,
\qquad
\operatorname{err}(\lambda)=K^{-1}\sum_{k=1}^K\operatorname{err}_k(\lambda),
\]
and choose $\lambda^*=\argmin_\lambda \operatorname{err}(\lambda)$. (See Notes pp.~22--24.)

Information criteria use the fitted log-likelihood and the selected model size $k=\#\{j:\hat\beta_j\ne0\}$:
\[
\operatorname{AIC}=-2\ell_n(\hat\beta)+2k,
\qquad
\operatorname{BIC}=-2\ell_n(\hat\beta)+k\log n.
\]
\red{Possible typo in source p.~26: the BIC bullet says ``AIC is minimized,'' but the formula and context indicate that BIC should be minimized.} Extended BIC adds stronger model-size penalization for high-dimensional model spaces. (See Notes p.~26.)

\subsubsection{Variable-selection quality}
Let $\beta^*$ be the true coefficient vector, $M=\{j:\beta_j^*\ne0\}$, and $\hat M=\{j:\hat\beta_j\ne0\}$. Variable-selection consistency means $\hat M=M$. (See Notes p.~28.)

The notes track mistakes by false positives and false negatives:
\[
\operatorname{FP}=\#\{j:\hat\beta_j\ne0,\ \beta_j^*=0\},
\qquad
\operatorname{FN}=\#\{j:\hat\beta_j=0,\ \beta_j^*\ne0\}.
\]
Sensitivity is $1-\operatorname{FN}/\abs{M}$, and specificity is $1-\operatorname{FP}/\abs{M^c}$. (See Notes pp.~28--31.)

\subsubsection{Representative examples}
\begin{itemize}
    \item The R / \texttt{glmnet} example simulates $N=50$, $p=100$, and five nonzero coefficients, then uses cross validation to choose $\lambda$ and reports false negatives, false positives, sensitivity, and specificity. It illustrates that Lasso can recover the true signals while still selecting extra noise variables. (See Notes pp.~29--31.)
    \item The solution-path discussion contrasts ridge and Lasso: ridge paths are smooth, while Lasso paths are piecewise linear and become sparse as $\lambda$ increases. (See Notes pp.~20--21.)
\end{itemize}

\subsection{Computing Lasso and certifying optimality}
\subsubsection{Coordinate gradient descent and soft-thresholding}
For the standardized linear-regression Lasso with objective (See Notes pp.~27, 32--40.)
\[
R(\beta_0,\beta)+\lambda p(\beta)
=\frac{1}{2N}\sum_{i=1}^N(y_i-\beta_0-x_i^\top\beta)^2
+\lambda\sum_{j=1}^p\abs{\beta_j},
\]
coordinate gradient descent updates one coordinate at a time. Holding all other coordinates fixed, define
\[
\tilde y_i^{(-j)}=\tilde\beta_0+\sum_{\ell\ne j}x_{i\ell}\tilde\beta_\ell.
\]
Then the one-coordinate update is
\[
\tilde\beta_j^{\mathrm{new}}
=S\left(\frac{1}{N}\sum_{i=1}^N x_{ij}(y_i-\tilde y_i^{(-j)}),\lambda\right),
\]
where the soft-thresholding operator is
\[
S(z,\lambda)=\operatorname{sign}(z)(\abs{z}-\lambda)_+.
\]
This is the computational expression of the $L_1$ penalty: take a negative-gradient step in the coordinate and then threshold toward zero. (See Notes pp.~32--34.)

The residual update avoids recomputing $\tilde y_i^{(-j)}$ from scratch:
\[
y_i-\tilde y_i^{(-j)}=r_i+x_{ij}\tilde\beta_j,
\qquad
\frac{1}{N}\sum_{i=1}^N x_{ij}(y_i-\tilde y_i^{(-j)})
=\frac{1}{N}\sum_{i=1}^N x_{ij}r_i+\tilde\beta_j.
\]
The naive full cycle costs $O(pN)$, and covariance updates can reduce repeated inner-product work once variables enter the model. (See Notes pp.~34--36.)

For weighted least squares,
\[
\frac{1}{2}\sum_{i=1}^N w_i(y_i-\beta_0-x_i^\top\beta)^2
+\lambda\sum_{j=1}^p\abs{\beta_j},
\]
the coordinate update becomes
\[
\tilde\beta_j^{\mathrm{new}}
=
\frac{
S\left(\sum_{i=1}^N w_i x_{ij}(y_i-\tilde y_i^{(-j)}),\lambda\right)
}{
\sum_{i=1}^N w_i x_{ij}^2
}.
\]
This weighted update is the bridge to penalized generalized linear models. (See Notes p.~37.)

\begin{mybox}
\textbf{How to solve the linear-regression Lasso by coordinate descent.}
\begin{enumerate}
    \item Standardize the predictors so that each coordinate update has the soft-thresholding form. (See Notes pp.~27, 32--33.)
    \item Cycle through coordinates, update each $\tilde\beta_j$ using $S(z,\lambda)$, and keep current residuals so that each update is cheap. (See Notes pp.~34--36.)
    \item Use the pathwise strategy: compute
    \[
    \lambda_{\max}
    =\max_{j\le p}\left|\frac{1}{N}\sum_{i=1}^N x_{ij}y_i\right|,
    \]
    because all coefficients stay zero when $\lambda\ge\lambda_{\max}$, then search a decreasing grid from $\lambda_{\max}$. (See Notes p.~39.)
    \item Since the linear-regression Lasso objective is convex, coordinate descent converges to the global minimizer. (See Notes p.~40.)
\end{enumerate}
\end{mybox}

\subsubsection{Convex analysis and KKT conditions}
For a convex function $f$, a vector $g$ is a subgradient at $x$ if
\[
f(y)\ge f(x)+g^\top(y-x)\qquad \text{for all }y.
\]
The set of all subgradients is the subdifferential $\partial f(x)$; if $f$ is differentiable and convex, then $\partial f(x)=\{\nabla f(x)\}$. (See Notes pp.~42--45.)

The minimizer condition for nonsmooth convex objectives is
\[
x^*\in\argmin_x f(x)
\quad\Longleftrightarrow\quad
0\in\partial f(x^*).
\]
This is the direct generalization of the differentiable first-order condition $\nabla f(x^*)=0$. (See Notes pp.~46--47.)

For the Lasso objective
\[
\frac12\norm{y-X\beta}_2^2+\lambda_n\norm{\beta}_1,
\]
the KKT conditions are, for each coordinate $j$,
\[
(X^\top X\beta)_j-X_j^\top y+\lambda_n\operatorname{sign}(\beta_j)=0
\quad \text{if }\beta_j\ne0,
\]
and
\[
\abs{(X^\top X\beta)_j-X_j^\top y}<\lambda_n
\quad \text{if }\beta_j=0.
\]
The strict inequality version is sufficient; replacing $<$ by $\le$ gives the necessary boundary version. Conceptually, the KKT conditions say
\[
0\in \nabla_\beta \frac12\norm{y-X\beta}_2^2+\lambda_n\partial\norm{\beta}_1.
\]
These conditions are the algebraic engine behind the Lasso selection-consistency proof. (See Notes pp.~52--53.)

\subsection{Non-asymptotic Lasso selection theory}
\subsubsection{Setup and assumptions}
The theoretical section studies (See Notes pp.~41, 48--49.)
\[
y=X\beta^*+\epsilon,
\qquad
\epsilon\sim N(0,\sigma^2I_n),
\qquad
\hat\beta
=\argmin_{\beta\in\R^p}
\left\{\frac12\norm{y-X\beta}_2^2+\lambda_n\norm{\beta}_1\right\}.
\]
Let $S=\{j:\beta_j^*\ne0\}$, $S^c=\{1,\ldots,p\}\setminus S$, $s=\abs{S}$, and $b_s=\min_{j\in S}\abs{\beta_j^*}$.

The main assumptions are:
\begin{itemize}
    \item (C1) $\norm{(X_S^\top X_S)^{-1}}_\infty=O(\sqrt{s}\,n^{-1})$, so the active-design Hessian is not too close to singular. (See Notes pp.~49--50.)
    \item (C2) $\norm{(X_{S^c}^\top X_S)(X_S^\top X_S)^{-1}}_\infty\le 1-\eta$ for some $0<\eta<1$. This is the irrepresentable condition: inactive covariates cannot be too well represented by active covariates. (See Notes pp.~49--51.)
    \item (C3) $b_s\ge c_0>0$, a minimum-signal-strength condition needed to avoid selecting zero when the true coefficient is too small. (See Notes pp.~49--51.)
\end{itemize}

\subsubsection{Theorem 1: non-asymptotic selection consistency}
\textbf{Theorem 1.}
If (C1)--(C3) hold and
\[
\lambda_n=C_1 n\left(\frac{\log p}{n}\right)^{1/2-\alpha}
\qquad\text{for }0<\alpha<1/2,
\]
then for some constant $C>\sqrt{2}$, with probability at least $1-p^{1-C^2/2}$, (See Notes p.~49.)
\[
\norm{\hat\beta_S-\beta_S^*}_\infty
\lesssim
\sqrt{s}\left(\frac{\log p}{n}\right)^{1/2-\alpha},
\qquad
\hat S=S.
\]
This is a non-asymptotic probability statement, but it implies $\norm{\hat\beta_S-\beta_S^*}_\infty=o_P(1)$ and model-selection consistency when the sparsity and dimension satisfy the displayed rate condition. In particular, $p$ may grow with $n$, even $p\gg n$, so long as $\sqrt{s}\{(\log p)/n\}^{1/2-\alpha}=o(1)$. (See Notes pp.~49--50.)

The theorem should be read as a statement about both signal and design. Sparsity alone is not enough: the active sub-design must be invertible enough, inactive variables must not be too correlated with active variables, and nonzero signals must clear the beta-min threshold. (See Notes pp.~50--51.)

\subsubsection{Proof skeleton}
The proof uses KKT conditions plus concentration. (See Notes pp.~54--61.)
First, define the event
\[
A=\left\{\norm{X^\top \epsilon}_\infty\le C\sigma\sqrt{n\log p}\right\}.
\]
By a Gaussian tail bound and the union bound,
\[
\prob(A)\ge 1-p^{1-C^2/2}.
\]
Conditioning on $A$, the proof first solves the active-set KKT equations inside an $\ell_\infty$ ball around $\beta_S^*$, using Miranda's theorem. It then extends the solution to $\hat\beta=(\hat\beta_S,0)^\top$ and verifies the inactive KKT inequalities by the irrepresentable condition and the fact that $\norm{X^\top\epsilon}_\infty=o(\lambda_n)$. (See Notes pp.~54--61.)

\begin{mybox}
\textbf{How to recognize a Lasso selection-consistency proof.}
\begin{enumerate}
    \item Write the KKT conditions separately on $S$ and $S^c$. (See Notes pp.~52--54.)
    \item Use a concentration bound to control $\norm{X^\top\epsilon}_\infty$ uniformly over $p$ covariates. (See Notes pp.~55--56.)
    \item Use the active-set Hessian condition to control $\norm{\hat\beta_S-\beta_S^*}_\infty$. (See Notes pp.~57--58.)
    \item Use the irrepresentable condition to force the inactive KKT inequalities and hence $\hat\beta_{S^c}=0$. (See Notes pp.~59--61.)
\end{enumerate}
\end{mybox}

\subsection{Penalized GLMs and generic optimization}
\subsubsection{High-dimensional generalized linear regression}
For canonical GLMs, the notes write the conditional density as (See Notes pp.~62--64.)
\[
f(y,X;\beta)
=\prod_{i=1}^N
c(y_i)\exp\left\{\frac{y_i\theta_i-b(\theta_i)}{\phi}\right\},
\qquad
(\theta_1,\ldots,\theta_N)^\top=X\beta.
\]
Up to an affine transformation, the log-likelihood is
\[
\ell(\beta)=y^\top X\beta-e^\top b(X\beta),
\]
and the $L_1$-regularized MLE solves
\[
\hat\beta^{\mathrm{lasso}}
=\argmin_\beta\{-\ell(\beta)+\lambda\norm{\beta}_1\}.
\]

For logistic regression,
\[
\log\frac{\prob(y=1\mid x)}{\prob(y=0\mid x)}
=\beta_0+x^\top\beta,
\]
and
\[
\ell(\beta_0,\beta)
=\sum_{i=1}^N
\{y_i(\beta_0+x_i^\top\beta)-\log(1+\exp(\beta_0+x_i^\top\beta))\}.
\]
The penalized logistic objective is $-\ell(\beta_0,\beta)+\lambda\sum_j\abs{\beta_j}$. (See Notes p.~64.)

\subsubsection{Local quadratic approximation and IRLS}
At the current iterate $(\tilde\beta_0,\tilde\beta)$, the logistic negative log-likelihood is approximated by a weighted least squares problem (See Notes pp.~65--67.)
\[
-\ell(\beta_0,\beta)
\approx
\frac12\sum_{i=1}^N
w_i(z_i-\beta_0-x_i^\top\beta)^2
+C(\tilde\beta_0,\tilde\beta),
\]
where
\[
\tilde p(x_i)=
\frac{\exp(\tilde\beta_0+x_i^\top\tilde\beta)}
{1+\exp(\tilde\beta_0+x_i^\top\tilde\beta)},
\quad
w_i=\tilde p(x_i)\{1-\tilde p(x_i)\},
\quad
z_i=\tilde\beta_0+x_i^\top\tilde\beta+
\frac{y_i-\tilde p(x_i)}{\tilde p(x_i)\{1-\tilde p(x_i)\}}.
\]
Thus penalized logistic regression alternates between updating the quadratic approximation and solving a penalized weighted least squares Lasso problem by coordinate descent. This is the penalized analogue of IRLS. (See Notes pp.~66--68.)

\subsubsection{MM and proximal gradient descent}
The MM algorithm constructs a majorizing function $\ell_Q(\beta\mid\tilde\beta^{(k)})$ satisfying
\[
\ell_Q(\beta\mid\tilde\beta^{(k)})\ge \ell(\beta),
\qquad
\ell_Q(\tilde\beta^{(k)}\mid\tilde\beta^{(k)})
=\ell(\tilde\beta^{(k)}).
\]
Minimizing the majorizer decreases the original objective because
\[
\ell(\tilde\beta^{(k)})
\le \ell(\tilde\beta^{(k-1)}).
\]
For twice continuously differentiable $\ell$, one can use a quadratic majorizer with curvature $c\ge \sup_\beta \lambda_{\max}\{\nabla^2\ell(\beta)\}$, which leads to a gradient step of size $1/c$. (See Notes pp.~69--71.)

For a nonsmooth objective
\[
\argmin_\beta\{\ell(\beta)+p(\beta)\},
\]
proximal gradient descent approximates the smooth part at $\tilde\beta^{(k)}$ and solves
\[
\argmin_\beta
\left\{
\frac{1}{2\eta}\norm{z-\beta}_2^2+p(\beta)
\right\},
\qquad
z=\tilde\beta^{(k)}-\eta\nabla\ell(\tilde\beta^{(k)}).
\]
The solution is the proximal operator $\operatorname{prox}_{\eta p}(z)$. (See Notes pp.~72--75.)

\textbf{Example 1.}
For Lasso, $p(\beta)=\lambda\norm{\beta}_1$, and the proximal operator is coordinatewise soft-thresholding:
\[
\operatorname{prox}_{\eta\lambda}(z)
=S_{\eta\lambda}(z)
=\operatorname{sign}(z)(\abs{z}-\eta\lambda)_+.
\]
\red{Possible typo in source p.~74: the displayed soft-thresholding formula omits $\abs{z}$, but p.~33 gives the correct $\operatorname{sign}(z)(\abs{z}-\lambda)_+$ form.} (See Notes pp.~33, 74.)

\begin{mybox}
\textbf{How to use proximal gradient descent.}
\begin{enumerate}
    \item Split the objective into a differentiable loss $\ell(\beta)$ and a possibly nonsmooth penalty $p(\beta)$. (See Notes pp.~72--73.)
    \item Take a gradient step on $\ell$: $z=\tilde\beta^{(k)}-\eta\nabla\ell(\tilde\beta^{(k)})$. (See Notes pp.~72--75.)
    \item Apply the proximal projection $\operatorname{prox}_{\eta p}(z)$. For $L_1$ penalties, this is soft-thresholding. (See Notes pp.~73--75.)
    \item Choose $\eta$ by backtracking line search: shrink $\eta$ by a factor $\alpha\in(0,1)$ until the quadratic majorization condition is met. (See Notes pp.~76--77.)
    \item Stop when the relative iterate change, such as $\norm{\tilde\beta^{(k+1)}-\tilde\beta^{(k)}}_2/\norm{\tilde\beta^{(k)}}_2$, is below a pre-specified tolerance. (See Notes p.~77.)
\end{enumerate}
\end{mybox}

\subsection{SCAD, adaptive Lasso, and oracle behavior}
\subsubsection{Penalty design requirements}
For the one-dimensional penalized problem (See Notes pp.~79--83.)
\[
\hat\theta=\argmin_\theta
\left\{\frac12(z-\theta)^2+p_\lambda(\abs{\theta})\right\},
\]
a good penalty should satisfy three goals:
\begin{itemize}
    \item \underline{Unbiasedness:} $p_\lambda'(\abs{\theta})=0$ for large $\abs{\theta}$, so large true signals are not unnecessarily shrunk. (See Notes pp.~80--82.)
    \item \underline{Sparsity:} $p_\lambda(\cdot)$ should be singular at the origin, allowing a thresholding rule that sets small estimates to zero. (See Notes pp.~80--82.)
    \item \underline{Continuity:} $\abs{\theta}+p_\lambda'(\abs{\theta})$ should attain its minimum at the origin, avoiding discontinuous prediction instability. (See Notes pp.~80--82.)
\end{itemize}
The $L_q$ penalties cannot satisfy all three: when $q>1$ they are not sparse, and when $q=1$ Lasso remains biased for large coefficients. Hard thresholding is sparse and less biased but discontinuous. (See Notes pp.~83--85.)

\subsubsection{SCAD penalty and computation}
Fan and Li's SCAD penalty has derivative, for $\theta>0$, (See Notes pp.~83--85.)
\[
p_\lambda'(\theta)
=
\lambda\left\{
I(\theta\le\lambda)
+
\frac{(a\lambda-\theta)_+}{(a-1)\lambda}I(\theta>\lambda)
\right\},
\qquad a>2.
\]
It is a quadratic spline with knots at $\lambda$ and $a\lambda$, and the recommended value is $a=3.7$.

For the one-dimensional problem, the SCAD thresholding operator is
\[
\hat\theta=
\begin{cases}
\operatorname{sign}(z)(\abs{z}-\lambda)_+, & \abs{z}\le 2\lambda,\\
\{(a-1)z-\operatorname{sign}(z)a\lambda\}/(a-2), & 2\lambda<\abs{z}\le a\lambda,\\
z, & \abs{z}>a\lambda.
\end{cases}
\]
Because SCAD is nonconvex, the full penalized likelihood problem can have multiple local optimizers; algorithms generally target a local minimizer / maximizer rather than a guaranteed global optimizer. (See Notes pp.~85--89.)

For
\[
\hat\beta^{\mathrm{scad}}
=\argmin_\beta
\left\{-L(\beta)+\sum_{j=1}^p p_\lambda(\abs{\beta_j})\right\},
\]
the notes use local quadratic approximation of $-L(\beta)$, solve the resulting weighted least squares problem by coordinate descent, and use the SCAD thresholding rule coordinatewise. The R package mentioned is \texttt{ncvreg}. (See Notes pp.~88--89.)

\subsubsection{Oracle properties of SCAD}
The source defines an oracle estimator as follows. Suppose $\beta_0=(\beta_{10}^\top,\beta_{20}^\top)^\top$ with $\beta_{20}=0$. An oracle estimator should satisfy $\hat\beta_2=0$ asymptotically and have $\hat\beta_1$ root-$n$ consistent for $\beta_{10}$. In words, it behaves as if one knew in advance which coefficients were zero. (See Notes p.~90.)

Let $V_i=(X_i,Y_i)$ be i.i.d.\ with density $f(V,\beta)$ and define
\[
Q(\beta)=L(\beta)-n\sum_{j=1}^p p_{\lambda_n}(\abs{\beta_j}).
\]

\textbf{Theorem 1 (SCAD local maximizer).}
Under regularity conditions (A)--(C), if
\[
\max\left\{\abs{p_{\lambda_n}''(\abs{\beta_{j0}})}:\beta_{j0}\ne0\right\}\to0,
\]
then there exists a local maximizer $\hat\beta$ of $Q(\beta)$ such that
\[
\norm{\hat\beta-\beta_0}_2
=O_P(n^{-1/2}+a_n),
\qquad
a_n=
\max\left\{\abs{p_{\lambda_n}'(\abs{\beta_{j0}})}:\beta_{j0}\ne0\right\}.
\]
The regularity conditions are the usual likelihood conditions: common support / identifiability, finite positive definite Fisher information at $\beta_0$, and controlled third derivatives. (See Notes pp.~91--95.)

The proof is a local-existence argument: show that with high probability $Q(\beta_0+\alpha_nu)<Q(\beta_0)$ on a sphere $\norm{u}_2=C$, where $\alpha_n=n^{-1/2}+a_n$. Then a local maximizer must lie inside that ball. (See Notes pp.~94--95.)

\textbf{Theorem 2 (oracle property).}
Let
\[
\Sigma=\operatorname{diag}\{p_{\lambda_n}''(\abs{\beta_{10}}),\ldots,
p_{\lambda_n}''(\abs{\beta_{s0}})\},
\]
and
\[
b=
\left(
p_{\lambda_n}'(\abs{\beta_{10}})\operatorname{sign}(\beta_{10}),
\ldots,
p_{\lambda_n}'(\abs{\beta_{s0}})\operatorname{sign}(\beta_{s0})
\right)^\top.
\]
If
\[
\liminf_{n\to\infty}\liminf_{\theta\to0+}
\frac{p_{\lambda_n}'(\theta)}{\lambda_n}>0,
\qquad
\lambda_n\to0,
\qquad
\sqrt n\,\lambda_n\to\infty,
\]
then with probability tending to one, the root-$n$ consistent local maximizer $\hat\beta=(\hat\beta_1^\top,\hat\beta_2^\top)^\top$ from Theorem 1 satisfies: (See Notes pp.~96--101.)
\begin{enumerate}
    \item[(a)] Sparsity: $\hat\beta_2=0$.
    \item[(b)] Asymptotic normality:
    \[
    \sqrt n\{I_1(\beta_{10})+\Sigma\}
    \left[
    \hat\beta_1-\beta_{10}
    +\{I_1(\beta_{10})+\Sigma\}^{-1}b
    \right]
    \overset{D}{\to}
    N(0,I_1(\beta_{10})),
    \]
    where $I_1(\beta_{10})=I_1(\beta_{10},0)$ is the Fisher information when $\beta_2=0$ is known.
\end{enumerate}

\textbf{Lemma 1.}
Under the conditions of Theorem 1, with probability tending to one, for any $\beta_1$ satisfying $\norm{\beta_1-\beta_{10}}_2=O_P(n^{-1/2})$ and any constant $C$,
\[
Q\left\{\begin{pmatrix}\beta_1\\0\end{pmatrix}\right\}
=
\max_{\norm{\beta_2}_2\le Cn^{-1/2}}
Q\left\{\begin{pmatrix}\beta_1\\\beta_2\end{pmatrix}\right\}.
\]
The lemma proves the sparsity part by showing that, near zero, the penalized objective is maximized at $\beta_2=0$ in the inactive coordinates. (See Notes pp.~97--100.)

\begin{mybox}
\textbf{How the SCAD oracle argument works.}
\begin{enumerate}
    \item Theorem 1 first establishes that a good local maximizer exists near the true parameter. (See Notes pp.~91--95.)
    \item Lemma 1 then shows that, locally, inactive coordinates are maximized at zero because the penalty derivative near zero dominates the random likelihood score. (See Notes pp.~96--100.)
    \item Theorem 2 combines sparsity with the likelihood expansion on the active coordinates to obtain the asymptotic normality statement. (See Notes pp.~100--102.)
    \item SCAD can satisfy both $\lambda_n\to0$ and $\sqrt n\lambda_n\to\infty$ while having $a_n=0$ for large true signals; ordinary Lasso has $a_n=\lambda_n$, which creates incompatible requirements for root-$n$ consistency and oracle sparsity. (See Notes p.~102.)
\end{enumerate}
\end{mybox}

\subsubsection{Adaptive Lasso}
Adaptive Lasso uses data-dependent weights to reduce the large-signal bias of ordinary Lasso: (See Notes pp.~103--106.)
\[
\hat\beta^{\mathrm{aLasso}}
=\argmin_\beta
\left\{
\sum_{i=1}^n(y_i-x_i^\top\beta)^2
+\lambda\sum_{j=1}^p \hat w_j\abs{\beta_j}
\right\},
\qquad
\hat w_j=\frac{1}{\abs{\hat\beta_j^{\mathrm{init}}}^{\gamma}},
\quad
\gamma>0.
\]
Large preliminary coefficients get smaller weights, while preliminary coefficients close to zero get large weights. This is the key mechanism behind simultaneous sparsity and reduced bias. (See Notes pp.~103--106.)

The source states that adaptive Lasso has oracle properties, referring to Theorem 2 of Zou (2006). Its practical advantages over SCAD are that the optimization problem is convex for fixed weights and it avoids the multiple-local-minimizer issue, but it requires an initial estimator; in high dimensions, least squares cannot serve as that initial estimator when $p\gg n$. (See Notes pp.~104--105.)

\subsection{Elastic net and grouping}
\subsubsection{Naive elastic net}
Lasso can be unstable with highly correlated variables because it tends to select one variable from a correlated group and ignore the others. Elastic net adds an $L_2$ penalty to encourage grouping: (See Notes pp.~107--108.)
\[
\hat\beta^{\mathrm{Naive\ ENet}}
=\argmin_\beta
\left\{
\norm{y-X\beta}_2^2
+\lambda_1\norm{\beta}_1
+\lambda_2\norm{\beta}_2^2
\right\}.
\]
The $L_1$ part creates sparsity, while the quadratic part stabilizes the path and encourages correlated predictors to enter together.

The constrained geometry is
\[
\min\norm{y-X\beta}_2^2
\quad\text{subject to}\quad
J(\beta)\le t,
\qquad
J(\beta)=\alpha\norm{\beta}_1+(1-\alpha)\norm{\beta}_2^2,
\]
where $\alpha=\lambda_1/(\lambda_1+\lambda_2)$. The $L_1$ corners give sparsity, and the curved edges from the $L_2$ part give grouping. (See Notes p.~108.)

\subsubsection{Theorem 1 of Zou and Hastie: grouping effect}
A method has the grouping effect if coefficients of highly correlated variables are close, and exactly identical predictors receive identical coefficients up to sign changes for negative correlation. (See Notes pp.~109--110.)

\textbf{Theorem 1 of Zou and Hastie (2005).}
Suppose $y$ is centered and predictors are standardized. Let $\hat\beta(\lambda_1,\lambda_2)$ be the naive elastic net estimate. If
\[
\hat\beta_i(\lambda_1,\lambda_2)\hat\beta_j(\lambda_1,\lambda_2)>0,
\]
define
\[
D_{\lambda_1,\lambda_2}(i,j)
=
\frac{1}{\norm{y}_1}
\abs{\hat\beta_i(\lambda_1,\lambda_2)-\hat\beta_j(\lambda_1,\lambda_2)}.
\]
Then
\[
D_{\lambda_1,\lambda_2}(i,j)
\le
\frac{1}{\lambda_2}\sqrt{2(1-\rho)},
\qquad
\rho=x_i^\top x_j.
\]
Thus when $\rho\approx1$, the coefficient-path distance is close to zero. This quantifies the grouping effect that ordinary Lasso lacks. (See Notes pp.~109--110.)

\subsubsection{Lemma 1 of Zou and Hastie: connection to Lasso}
\textbf{Lemma 1 of Zou and Hastie (2005).}
Given $(y,X)$ and $(\lambda_1,\lambda_2)$, define the augmented data set
\[
X^*=(1+\lambda_2)^{-1/2}
\begin{pmatrix}
X\\
\sqrt{\lambda_2}I
\end{pmatrix},
\qquad
y^*=
\begin{pmatrix}
y\\
0
\end{pmatrix}.
\]
Let
\[
\gamma=\frac{\lambda_1}{\sqrt{1+\lambda_2}},
\qquad
\beta^*=\sqrt{1+\lambda_2}\,\beta.
\]
Then the naive elastic net criterion can be written as a Lasso problem
\[
\hat\beta^*
=\argmin_{\beta^*}
\left\{
\norm{y^*-X^*\beta^*}_2^2
+\gamma\norm{\beta^*}_1
\right\},
\]
and
\[
\hat\beta^{\mathrm{Naive\ ENet}}
=\frac{1}{\sqrt{1+\lambda_2}}\hat\beta^*.
\]
This is why elastic net can be computed through Lasso algorithms such as coordinate descent and \texttt{glmnet}. (See Notes pp.~111--113.)

The rescaled elastic net estimator is
\[
\hat\beta^{\mathrm{ENet}}
=(1+\lambda_2)\hat\beta^{\mathrm{Naive\ ENet}}.
\]
The rescaling is intended to keep the grouping effect while correcting some of the extra shrinkage from the quadratic penalty. Another interpretation is that elastic net replaces the sample covariance $\hat\Sigma=X^\top X$ by a regularized version
\[
\hat\Sigma_{\lambda_2}
=(1-\gamma)\hat\Sigma+\gamma I,
\qquad
\gamma=\frac{\lambda_2}{1+\lambda_2},
\]
before applying $L_1$ regularization. (See Notes p.~112.)

\begin{mybox}
\textbf{Final conceptual summary.}
\begin{itemize}
    \item \underline{Ridge:} $L_2$ regularization fixes ill-conditioning and improves prediction stability, but it is biased and does not select variables. (See Notes pp.~4--15.)
    \item \underline{Lasso:} $L_1$ regularization is the convex route to sparsity. Its computation is built around soft-thresholding, and its theory is built around KKT conditions plus concentration. (See Notes pp.~16--61.)
    \item \underline{GLM Lasso:} non-Gaussian likelihoods are handled by quadratic approximations or generic proximal-gradient logic; the nonsmooth penalty is handled by the proximal map. (See Notes pp.~62--78.)
    \item \underline{SCAD / adaptive Lasso:} these methods are designed to reduce Lasso's bias for large signals and can have oracle properties under tuning and regularity conditions, but SCAD introduces nonconvex local-optimizer issues. (See Notes pp.~79--106.)
    \item \underline{Elastic net:} adding an $L_2$ penalty to Lasso gives grouping behavior for correlated predictors and reduces to an augmented-data Lasso problem computationally. (See Notes pp.~107--113.)
\end{itemize}
\end{mybox}
